\chapter{OPERATOR KOMPAK}\label{chap_cpt_ops}

\section{Definisi dan Sifat-Sifat Dasar}

Kita akan melanjutkan pembahasan dari akhir Bab 5 dengan mempelajari kelas-kelas operator.
Konteks kita akan sedikit lebih umum; dalam hal operator kompak, kita akan
meninjau operator pada ruang Banach.  Pertama, mari kita ingat kembali beberapa definisi dan
fakta dasar dari kalkulus lanjut (atau mata kuliah pengantar analisis).

\begin{defn} Suatu subhimpunan $A$ dari ruang metrik $M$ disebut
 \index{total!keterbatasan}%
 \index{terbatas!total}%
\df{terbatas total} jika untuk setiap $\epsilon > 0$ terdapat suatu subhimpunan hingga $F$ dari $A$
sedemikian sehingga untuk setiap $a \in A$ terdapat titik $x \in F$ yang memenuhi $d(x,a) < \epsilon$.
Untuk ruang linear bernorma, definisi ini dapat dirumuskan kembali: Suatu subhimpunan $A$ dari ruang linear
bernorma $V$ disebut \df{terbatas total} jika untuk setiap bola terbuka $B$ di sekitar nol dalam $V$ terdapat
suatu subhimpunan hingga $F$ dari $A$ sedemikian sehingga $A \subseteq F + B$.  Hal ini memiliki parafrasa yang kurang lebih baku:
Suatu ruang bersifat terbatas total jika dapat diawasi oleh sejumlah hingga polisi dengan jarak pandang
sependek apa pun.  (Sebagian penulis menyebut sifat ini \emph{prakompak}.)
\end{defn}

\begin{prop} Suatu ruang metrik bersifat kompak jika dan hanya jika ruang itu lengkap dan terbatas total.
\end{prop}

\begin{prop} Suatu ruang metrik bersifat kompak jika dan hanya jika ruang itu lengkap dan terbatas total.
\end{prop}

\begin{defn} Suatu subhimpunan $A$ dari ruang topologis $X$ disebut
 \index{kompak!relatif}%
 \index{relatif kompak}%
\df{kompak relatif} jika tutupannya kompak.
\end{defn}

\begin{prop} Dalam ruang metrik, setiap himpunan kompak relatif adalah terbatas total.
\end{prop}

Dalam ruang metrik lengkap, kebalikannya berlaku.

\begin{prop} Dalam ruang metrik lengkap, setiap himpunan terbatas total adalah kompak relatif.
\end{prop}

\begin{prop} Dalam ruang metrik, setiap himpunan terbatas total adalah separabel.
\end{prop}

\begin{thm}[Teorema Arzel\`a-Ascoli]\label{AAThm} Misalkan $X$ ruang Hausdorff kompak. Suatu subhimpunan $\fml F$
dari $\fml C(X)$ terbatas total apabila subhimpunan itu terbatas (secara seragam) dan ekuikontinu.
\end{thm}

Berikutnya kita akan membandingkan topologi \emph{lemah} dan \emph{kuat} pada ruang linear bernorma.
Istilah
 \index{lemah!\emph{vs.} topologi kuat pada ruang linear bernorma}%
 \index{topologi!lemah!\emph{vs.} kuat pada ruang linear bernorma}%
 \index{kuat!\emph{vs.} topologi lemah pada ruang linear bernorma}%
 \index{konvensi!pada ruang linear bernorma \emph{topologi kuat} merujuk pada topologi norma}%
\emph{topologi lemah} akan merujuk pada topologi lemah biasa yang ditentukan oleh fungsional linear terbatas
pada ruang tersebut sebagaimana didefinisikan dalam~\ref{C067434}, sedangkan \emph{topologi kuat} akan merujuk pada topologi norma pada ruang tersebut.
Namun, perlu diingat bahwa kata ``kuat'' dalam konteks ruang linear bernorma, jika ditilik secara ketat,
sesungguhnya mubazir.  Kata itu digunakan semata-mata sebagai penegasan.  Misalnya, banyak penulis akan menyatakan proposisi berikut sebagai:
\emph{Setiap barisan yang konvergen secara lemah dalam ruang linear bernorma adalah terbatas.}

\begin{prop} Setiap barisan yang konvergen secara lemah dalam ruang linear bernorma adalah terbatas secara kuat.
\end{prop}

\begin{proof}[\emph{Petunjuk untuk bukti}] Gunakan \emph{prinsip keterbatasan seragam}~\ref{thm_PUB}. \ns
\end{proof}

\begin{defn} Kita mengatakan bahwa suatu pemetaan linear $T\colon V \sto W$ bersifat
 \index{lemah!kekontinuan}%
 \index{kontinu!secara lemah}%
 \index{linear!pemetaan!kekontinuan lemah suatu}%
\df{kontinu secara lemah} jika pemetaan itu membawa jaring yang konvergen secara lemah ke jaring yang konvergen secara lemah.
\end{defn}

\begin{prop} Setiap pemetaan linear yang kontinu (secara kuat) antara ruang-ruang linear bernorma bersifat kontinu secara lemah.
\end{prop}

\begin{defn} Suatu pemetaan linear $T\colon V \sto W$ antara ruang-ruang linear bernorma disebut
 \index{kompak!pemetaan linear}%
 \index{linear!pemetaan!kompak}%
 \index{operator!kompak}%
\df{kompak} jika pemetaan itu membawa himpunan terbatas ke himpunan kompak relatif. Secara ekuivalen, $T$ bersifat kompak jika pemetaan itu
memetakan bola satuan tertutup dalam $V$ ke suatu himpunan kompak relatif dalam~$W$.   Keluarga semua pemetaan linear kompak
dari $V$ ke $W$ dinotasikan
 \index{K@$\ofml K(V,W)$!keluarga pemetaan linear kompak}%
 \index{K@$\ofml K(V)$!keluarga operator kompak}%
dengan~$\ofml K(V,W)$; dan kita tulis $\ofml K(V)$ untuk~$\ofml K(V,V)$.
\end{defn}

\begin{prop} Setiap pemetaan linear kompak bersifat terbatas.
\end{prop}

\begin{prop} Suatu pemetaan linear $T\colon V \sto W$ antara ruang-ruang linear bernorma bersifat kompak jika dan hanya jika pemetaan itu membawa
setiap barisan terbatas dalam $V$ ke suatu barisan dalam $W$ yang memiliki subbarisan konvergen.
\end{prop}

\begin{exam} Setiap operator berperingkat hingga pada ruang Banach bersifat kompak.
\end{exam}

\begin{exam} Setiap fungsional linear terbatas pada ruang Banach $B$ merupakan pemetaan linear kompak dari $B$ ke~$\K$.
\end{exam}

\begin{exam} Misalkan $C$ ruang Banach $\fml C([0,1])$ dan $k \in \fml C([0,1] \times [0,1])$.  Untuk setiap $f \in C$
 \index{integral!operator}%
 \index{operator!integral}%
definisikan fungsi $Kf$ pada $[0,1]$ dengan
   \[ Kf(x) = \int_0^1 k(x,y)f(y)\,dy\,. \]
Maka operator integral $K$ merupakan operator kompak pada~$C$.
\end{exam}

\begin{proof}[\emph{Petunjuk untuk bukti}] Gunakan \emph{teorema Arzel\`a-Ascoli}~\ref{AAThm}.  \ns
\end{proof}

\begin{exam}\label{X_square_int_kernel} Misalkan $H$ ruang Hilbert $\lfs2{[0,1]}$ yang terdiri atas (kelas-kelas ekuivalensi) fungsi yang
terintegralkan kuadrat pada $[0,1]$ terhadap ukuran Lebesgue dan
 \index{integral!operator}%
 \index{operator!integral}%
$k$ suatu fungsi yang terintegralkan kuadrat
pada $[0,1] \times [0,1]$.  Untuk setiap $f \in H$ definisikan fungsi $Kf$ pada $[0,1]$ dengan
   \[ Kf(x) = \int_0^1 k(x,y)f(y)\,dy\,. \]
Maka operator integral $K$ merupakan operator kompak pada~$H$.
\end{exam}

\begin{prop} Suatu operator pada ruang Hilbert bersifat kompak jika dan hanya jika operator itu memetakan bola satuan tertutup dalam ruang tersebut
ke suatu himpunan kompak.
\end{prop}

\begin{cor}\label{cor_id_op_not_cpt} Operator identitas pada ruang Hilbert berdimensi tak hingga tidak kompak.
(Lihat contoh~\ref{X_l2ball_not_cpt}.)
\end{cor}

\begin{exam} Jika $B$ ruang Banach, keluarga $\ofml B(B)$ dari operator pada $B$ merupakan
 \index{Banach!aljabar!$\ofml B(B)$ sebagai suatu}%
 \index{B@$\ofml B(V)$!sebagai aljabar Banach beridentitas}%
aljabar Banach beridentitas dan
 \index{K@$\ofml K(V)$!sebagai ideal tertutup sejati}%
keluarga $\ofml K(B)$ dari operator kompak pada $B$ merupakan ideal tertutup dalam~$\ofml B(B)$.
Jika $B$ berdimensi tak hingga, ideal $\ofml K(B)$ adalah sejati.
\end{exam}

\begin{prop} Misalkan $T\colon B \sto C$ suatu pemetaan linear terbatas antara ruang-ruang Banach. Maka $T$ bersifat kompak
jika dan hanya jika $T^*$ bersifat kompak.
\end{prop}

\begin{defn}\label{0013} Suatu
 \index{C*@aljabar-$C^*$}%
 \index{aljabar!$C^*$-}%
\df{aljabar-$C^\ast$} adalah aljabar Banach $A$ dengan involusi yang memenuhi
   \[ \norm{a^*a} = \norm{a}^2 \]
untuk setiap $a \in A$.  Sifat norma ini biasanya disebut
 \index{C*@syarat-$C^*$}%
\emph{syarat-$C^*$}. Norma aljabar yang memenuhi syarat ini adalah
 \index{C*@norma-$C^*$}%
 \index{norma!$C^*$-}%
\df{norma-$C^*$}.  Suatu
 \index{C*@subaljabar-$C^*$}%
 \index{subaljabar!$C^*$-}%
\df{subaljabar-$C^*$} dari aljabar-$C^*$ $A$ adalah subaljabar-$*\,$ tertutup dari~$A$.

Kita menyatakan kategori aljabar-$C^*$ dan homomorfisme-$*\,$ dengan $\cat{CSA}$.
 \index{kategori!$\cat{CSA}$ sebagai suatu}%
 \index{CSA@$\cat{CSA}$!kategori}%
Sepintas mungkin tampak aneh bahwa dalam kategori ini, yang objek-objeknya memiliki sifat aljabar sekaligus topologis,
morfismenya hanya disyaratkan mempertahankan struktur aljabar. Namun, ternyata
setiap morfisme dalam $\cat{CSA}$ secara otomatis kontinu---bahkan, kontraktif (lihat
proposisi~\futurexref{12.3.16}{00152171})---dan
 \index{morfisme bijektif!dapat dibalik!dalam $\cat{CSA}$}%
setiap morfisme injektif merupakan isometri (lihat proposisi~\futurexref{12.3.17}{00152181}).  Jelas bahwa
(seperti pada definisi~\ref{00109}) setiap morfisme bijektif dalam kategori ini merupakan isomorfisme. Dengan demikian, setiap homomorfisme-$*\,$ bijektif
merupakan isomorfisme-$*\,$ isometrik.
\end{defn}

\begin{exam} Ruang vektor $\C$ dari bilangan kompleks, dengan perkalian bilangan kompleks biasa
 \index{C*@aljabar-$C^*$!$\C$ sebagai suatu}%
 \index{C@$\C$!sebagai aljabar-$C^*$}%
dan konjugasi kompleks $z \mapsto \conj z$ sebagai involusi, merupakan aljabar-$C^*$
komutatif beridentitas.
\end{exam}

\begin{exam} Jika $X$ suatu ruang Hausdorff kompak, aljabar $\fml C(X)$ dari fungsi kontinu
 \index{C*@aljabar-$C^*$!$\fml C(X)$ sebagai suatu}%
 \index{C@$\fml C(X)$!sebagai aljabar-$C^*$}%
bernilai kompleks pada $X$ merupakan aljabar-$C^*$ komutatif beridentitas apabila involusinya
diambil sebagai konjugasi kompleks.
\end{exam}

\begin{exam} Jika $X$ suatu ruang Hausdorff kompak lokal, aljabar Banach
 \index{C*@aljabar-$C^*$!$\fml C_0(X)$ sebagai suatu}%
 \index{C@$\fml C_0(X)$!sebagai aljabar-$C^*$}%
$\fml C_0(X) = \fml C_0(X,\C)$ dari fungsi kontinu bernilai kompleks pada $X$ yang lenyap
di tak hingga merupakan aljabar-$C^*$ komutatif (tidak mesti beridentitas) apabila involusinya
diambil sebagai konjugasi kompleks.
\end{exam}

\begin{exam} Jika $(X,\mu)$ suatu ruang ukur, aljabar $\lfs\infty(X,\mu)$ dari
 \index{C*@aljabar-$C^*$!$L_\infty(S)$ sebagai suatu}%
 \index{L@$\lfs\infty(S)$!sebagai aljabar-$C^*$}%
fungsi terukur bernilai kompleks yang terbatas secara esensial pada $X$ (sekali lagi dengan konjugasi
kompleks sebagai involusi) merupakan aljabar-$C^*$.  (Secara teknis, tentu saja, anggota
$\lfs\infty(X,\mu)$ adalah kelas-kelas ekuivalensi fungsi yang berbeda pada himpunan-himpunan berukuran nol.)
\end{exam}

\begin{exam}\label{001305fa} Aljabar $\ofml B(H)$ dari operator linear terbatas pada ruang Hilbert
 \index{C*@aljabar-$C^*$!$\ofml B(H)$ sebagai suatu}%
 \index{B@$\ofml B(V)$!sebagai aljabar-$C^*$}%
$H$ merupakan aljabar-$C^*$ beridentitas.  Keluarga $\ofml K(H)$ dari operator kompak merupakan ideal-$*\,$ tertutup
dalam aljabar tersebut.
\end{exam}

\begin{exam}\label{001306} Sebagai kasus khusus dari contoh sebelumnya, perhatikan bahwa himpunan
 \index{C*@aljabar-$C^*$!M@$\mathbf M_n$ sebagai suatu}%
 \index{Ma@$\mathbf M_n = \M n\C$!sebagai aljabar-$C^*$}%
$\mathbf M_n$ dari matriks bilangan kompleks berukuran $n \times n$ merupakan aljabar-$C^*$ beridentitas.
\end{exam}

\begin{prop}\label{prop_K_is_clo_FR} Dalam ruang Hilbert $H$, ideal $\ofml K(H)$ dari operator kompak merupakan tutupan
ideal $\ofml{FR}(H)$ dari operator berperingkat hingga.
\end{prop}

Selama kira-kira satu dasawarsa sebelum Perang Dunia Kedua, sejumlah matematikawan Polandia terkemuka berkumpul secara teratur di kedai
kopi di Lwow untuk membahas hasil dan mengajukan masalah.  Mula-mula mereka bertemu di \emph{Caf\'e Roma}, kemudian di
\emph{Scottish Caf\'e}.  Pada 1935 Stefan Banach menyediakan buku catatan berjilid tempat mereka dapat menuliskan masalah-masalah tersebut.
Dari sinilah \emph{Scottish Book}~\cite{Mauldin:1981} yang kini terkenal bermula.  Pada 6 November 1936, Stanislaw Mazur
menuliskan Masalah 153, yang dikenal sebagai \emph{masalah aproksimasi}; pada intinya masalah itu menanyakan apakah proposisi sebelumnya
berlaku untuk ruang Banach.  Dapatkah setiap operator kompak pada ruang Banach diaproksimasi dalam norma oleh operator berperingkat hingga?
Sebagai hadiah bagi pemecah masalah ini, Mazur menawarkan seekor angsa hidup (tergolong salah satu hadiah paling murah hati yang dijanjikan dalam
\emph{Scottish Book}).  Masalah tersebut tetap terbuka selama beberapa dasawarsa, hingga pada 1972 Per Enflo memberikan jawaban negatif,
setelah berhasil menemukan ruang Banach refleksif separabel yang di dalamnya aproksimasi semacam itu gagal.  Dalam sebuah upacara yang diadakan
tahun itu di Warsawa, Enflo menerima angsa hidupnya dari Mazur sendiri.  Hasil tersebut diterbitkan dalam Acta
Mathematica~\cite{Enflo:1973} pada~1973.

\begin{exam} Operator diagonal $\diag(1,\frac12,\frac13,\dots)$ merupakan operator kompak pada ruang Hilbert~$l_2$.
\end{exam}











\section{Isometri Parsial}
\begin{defn}  Suatu elemen $v$ dari aljabar-$C^*$ $A$ disebut
 \index{parsial!isometri}%
\df{isometri parsial} jika $v^*v$ merupakan proyeksi dalam~$A$.  (Karena $v^*v$ selalu
swaadjoin, cukup disyaratkan bahwa $v^*v$ idempoten.)  Elemen $v^*v$ adalah
\df{proyeksi}
 \index{awal!proyeksi}%
 \index{proyeksi!awal}%
 \index{parsial!isometri!proyeksi awal dari suatu}%
\df{awal} (atau
 \index{tumpuan!proyeksi}%
 \index{proyeksi!tumpuan}%
 \index{parsial!isometri!proyeksi tumpuan dari suatu}%
\df{tumpuan}) dari~$v$, sedangkan $vv^*$ adalah \df{proyeksi}
 \index{akhir!proyeksi}%
 \index{proyeksi!akhir}%
 \index{parsial!isometri!proyeksi akhir dari suatu}%
\df{akhir} (atau
 \index{jangkauan!proyeksi}%
 \index{proyeksi!jangkauan}%
 \index{parsial!isometri!proyeksi jangkauan dari suatu}%
\df{jangkauan}) dari~$v$.  (Sebagai konsekuensi langsung dari proposisi berikutnya,
jika $v$ suatu isometri parsial, maka $vv^*$ memang merupakan proyeksi.)
\end{defn}

\begin{prop}\label{pi0111} Jika $v$ suatu isometri parsial dalam aljabar-$C^*$, maka
 \begin{enumerate}[label=\rm{(\alph*)}]
  \item $vv^*v = v$\,; dan
  \item $v^*$ merupakan isometri parsial.
 \end{enumerate}
\end{prop}

\begin{proof}[\emph{Petunjuk untuk bukti}] Misalkan $z = v - vv^*v$ dan tinjau $z^*z$.  \ns
\end{proof}

\begin{prop}\label{pi0114} Misalkan $v$ suatu isometri parsial dalam aljabar-$C^*$~$A$.  Maka proyeksi awalnya
$p$ adalah proyeksi terkecil (terhadap urutan parsial~$\preceq$
pada $\fml P(A)$\,) sedemikian sehingga $vp = v$, dan proyeksi akhirnya $q$ adalah proyeksi
terkecil sedemikian sehingga $qv = v$.
\end{prop}

\begin{prop}\label{pi0117} Jika $V$ suatu isometri parsial pada ruang Hilbert $H$ (yakni, jika $V$ suatu isometri parsial
dalam aljabar-$C^*$~$\ofml B(H)$\,), maka proyeksi awal $V^*V$ adalah
proyeksi dari $H$ pada $(\ker V)^\perp$, dan proyeksi akhir $VV^*$ adalah proyeksi
dari $H$ pada~$\ran V$.
\end{prop}

Berdasarkan hasil sebelumnya, $(\ker V)^\perp$ disebut
 \index{awal!ruang}%
 \index{ruang!awal}%
 \index{ruang!tumpuan}%
 \index{tumpuan!ruang}%
\df{ruang} \df{awal} (atau \df{tumpuan}) dari~$V$, sedangkan $\ran V$ kadang-kadang disebut
 \index{akhir!ruang}%
 \index{ruang!akhir}%
\df{ruang akhir} dari~$V$.

\begin{prop}\label{pi0121} Suatu operator pada ruang Hilbert merupakan isometri parsial jika dan hanya jika operator tersebut merupakan
isometri pada komplemen ortogonal dari kernelnya.
\end{prop}





\begin{thm}[Dekomposisi Polar]\label{0022287} Jika $T$ suatu operator pada ruang
 \index{dekomposisi polar}%
 \index{dekomposisi!polar}%
Hilbert $H$, maka terdapat isometri parsial $V$ pada $H$ sedemikian sehingga
 \begin{enumerate}
    \item[(i)] ruang awal $V$ adalah $(\ker T)^\perp$,
    \item[(ii)] ruang akhir $V$ adalah $\clo{\ran T}$, dan
    \item[(iii)] $T = V\abs T$.
 \end{enumerate}
Dekomposisi $T$ ini unik dalam arti bahwa jika $T = V_0P$ dengan $V_0$ suatu
isometri parsial dan $P$ suatu operator positif sedemikian sehingga $\ker V_0 = \ker P$, maka $P =
\abs T$ dan $V_0 = V$.
\end{thm}

\begin{proof} Lihat \cite{Conway:1990}, VIII.3.11;\allowbreak \cite{Davidson:1996}, teorema I.8.1;
\allowbreak \cite{KadisonR:1983}, teorema 6.1.2;\allowbreak \cite{Murphy:1990}, teorema 2.3.4;
atau~\cite{Pedersen:1995}, teorema 3.2.17. \ns
\end{proof}

\begin{cor} Jika $T$ adalah suatu
  \index{dekomposisi polar!dari operator invertibel}%
 \index{operator!invertibel!dekomposisi polar dari}%
 \index{invertibel!operator!dekomposisi polar dari}%
operator invertibel pada suatu ruang Hilbert, maka isometri parsial dalam dekomposisi polar
(lihat~\ref{0022287}) bersifat uniter.
\end{cor}

\begin{prop} Jika $T$ adalah suatu
 \index{dekomposisi polar!dari operator normal}%
 \index{operator!normal!dekomposisi polar dari}%
 \index{normal!operator!dekomposisi polar dari}%
operator normal pada suatu ruang Hilbert $H$, maka terdapat operator uniter $U$ pada $H$
yang berkomutasi dengan $T$ dan memenuhi $T = U\abs T$.
\end{prop}

\begin{proof} Lihat \cite{Pedersen:1995}, proposisi 3.2.20.\ns \end{proof}

\begin{exer} Misalkan $(S, \fml A,\mu)$ suatu ruang ukur $\sigma$-hingga dan $\phi \in
L_\infty(\mu)$. Tentukan dekomposisi polar dari
 \index{operator perkalian!dekomposisi polar dari}%
 \index{operator!perkalian!dekomposisi polar dari}%
 \index{dekomposisi polar!dari operator perkalian}%
operator perkalian $M_\phi$ (lihat contoh~\ref{C063527}).
\end{exer}









\section{Operator Kelas Jejak}

Mari kita mulai dengan meninjau beberapa sifat standar \emph{jejak} suatu matriks $n \times n$.

\begin{defn}\label{tr0113} Misalkan $a = [a_{ij}] \in \mathbf M_n$. Definisikan $\tr a$, yaitu
 \index{jejak!dari suatu matriks}%
 \index{matriks!jejak dari suatu}%
 \index{tr@$\tr$ (jejak)}%
\df{jejak} dari $a$, dengan
  \[ \tr a = \sum_{i = 1}^n a_{ii}. \]
\end{defn}

\begin{prop}\label{tr0117} Fungsi $\tr\colon \mathbf M_n \sto \C$ bersifat linear.
\end{prop}

\begin{prop}\label{tr0121} Jika $a$, $b \in \mathbf M_n$, maka $\tr (ab) = \tr (ba)$.
\end{prop}

Jejak bersifat invarian terhadap keserupaan. Dua matriks $n \times n$, yaitu $a$ dan $b$, disebut
 \index{serupa!matriks}%
 \index{matriks!serupa}%
\df{serupa} jika terdapat matriks invertibel $s$ sedemikian sehingga $b = sas^{-1}$.

\begin{prop}\label{tr0124} Matriks-matriks dalam $\mathbf M_n$ yang serupa mempunyai jejak yang sama.
\end{prop}

\begin{defn}\label{tr0211} Misalkan $H$ suatu ruang Hilbert separabel, $(e_n)$ suatu basis ortonormal bagi~$H$, dan $T$
suatu operator positif pada~$H$. Definisikan
 \index{jejak}%
 \index{tr@$\tr$ (jejak)}%
\df{jejak} dari~$T$~dengan
    \[ \tr T:=\sum_{k=1}^\infty \langle Te_k,e_k \rangle. \]
(Jejak tidak harus berhingga.)
\end{defn}

\begin{prop} Jika $S$ dan $T$ merupakan operator positif pada suatu ruang Hilbert separabel dan $\alpha \ge 0$, maka
   \[  \tr(S + T) = \tr S + \tr T \]
dan
   \[ \tr(\alpha T) = \alpha \tr T. \]
\end{prop}

\begin{prop} Jika $T$ suatu operator pada ruang Hilbert separabel, maka $\tr(T^*T) = \tr(TT^*)$.
\end{prop}

Untuk menunjukkan bahwa definisi jejak operator positif yang diberikan di atas tidak bergantung pada basis ortonormal yang dipilih bagi
ruang Hilbert tersebut, kita memerlukan suatu hasil (yang diberikan berikut ini) yang bergantung pada fakta yang muncul kemudian dalam catatan ini: setiap operator positif pada ruang Hilbert
memiliki akar kuadrat (positif) (lihat contoh~\futurexref{11.5.7}{X_sqroot_op}).

\begin{prop}\label{tr0221} Pada keluarga operator positif di suatu ruang Hilbert separabel $H$, jejak bersifat invarian terhadap ekuivalensi uniter;
yakni, jika $T$ operator positif dan $U$ uniter, maka $\tr T = \tr(UTU^*)$.
\end{prop}

\begin{proof}[\emph{Petunjuk untuk bukti}] Dalam proposisi sebelumnya, ambil $S = UT^{\frac12}$.    \ns
\end{proof}

\begin{prop}\label{tr0224} Definisi \emph{jejak}~\ref{tr0211} dari suatu operator positif pada ruang Hilbert separabel tidak bergantung
pada pemilihan basis ortonormal bagi ruang tersebut.
\end{prop}

\begin{proof}[\emph{Petunjuk untuk bukti}] Misalkan $(e^k)$ dan $(f^k)$ basis-basis ortonormal bagi ruang tersebut dan $T$ suatu operator positif
pada ruang itu. Ambil $U$ sebagai operator uniter unik yang memenuhi $e^k = Uf^k$ untuk setiap $k \in \N$.  \ns
\end{proof}

\begin{defn} Misalkan $V$ suatu ruang vektor. Suatu subhimpunan $C$ dari $V$ disebut
 \index{kerucut}%
\df{kerucut} jika $\alpha C \subseteq C$ untuk setiap $\alpha \ge 0$. Suatu kerucut $C$ dalam $V$ disebut
 \index{kerucut!proper}%
 \index{proper!kerucut}%
\df{proper} jika $C \cap (-C) = \{\vc 0\}$.
\end{defn}

 \begin{prop} Suatu kerucut $C$ dalam ruang vektor bersifat konveks jika dan hanya jika $C + C \subseteq C$.
\end{prop}

\begin{exam} Pada suatu ruang Hilbert separabel $H$, keluarga semua operator positif dengan jejak hingga membentuk kerucut konveks proper dalam ruang
vektor~$\ofml B(H)$.
\end{exam}

\begin{defn} Dalam ruang Hilbert separabel $H$, subruang linear dari $\ofml B(H)$ yang direntang oleh operator-operator positif dengan jejak hingga
adalah keluarga $\ofml T(H)$ dari operator
 \index{jejak!kelas}%
 \index{operator!kelas jejak}%
 \index{T@$\ofml T(H)$ (operator kelas jejak pada~$H$)}%
\df{kelas jejak} pada~$H$.
\end{defn}

\begin{prop} Misalkan $H$ ruang Hilbert separabel dengan basis ortonormal~$(e_n)$. Jika $T$ operator kelas jejak pada $H$, maka
   \[ \tr T = \sum_{k=1}^\infty \langle Te_k,e_k \rangle \]
dan deret di ruas kanan konvergen mutlak.
\end{prop}
















\section{Operator Hilbert--Schmidt}

\begin{defn} Suatu operator $A$ pada ruang Hilbert separabel $H$ adalah operator
 \index{Hilbert!--Schmidt, operator}%
 \index{operator!Hilbert--Schmidt}%
\df{Hilbert--Schmidt} jika $A^*A$ termasuk kelas jejak; yakni, jika $\sum_{k=1}^\infty \norm{Ae_k}^2 < \infty$,
dengan $(e_k)$ suatu basis ortonormal untuk~$H$. Keluarga semua operator Hilbert--Schmidt pada $H$
 \index{HS@$\ofml{HS}(H)$ (operator Hilbert--Schmidt pada~$H$)}%
dinyatakan dengan~$\ofml{HS}(H)$.
\end{defn}

\begin{prop} Jika $H$ ruang Hilbert separabel, maka keluarga operator Hilbert--Schmidt pada $H$ merupakan ideal
dua sisi dalam~$\ofml B(H)$.
\end{prop}

\begin{exam} Keluarga operator Hilbert--Schmidt pada suatu ruang Hilbert $H$ dapat dijadikan ruang hasil kali dalam.
\end{exam}

\begin{proof}[\emph{Petunjuk untuk bukti}] Polarisasi. Jika $A$, $B \in \ofml{HS}(H)$,
tinjau $\sum_{k=0}^3 i^k \bigl(A + i^kB\bigr)^*(A + i^kB)$. \ns
\end{proof}

\begin{exam} Hasil kali dalam yang didefinisikan dalam contoh sebelumnya menjadikan $\ofml{HS}(H)$ suatu ruang Hilbert.
\end{exam}

\begin{prop} Pada ruang Hilbert separabel, setiap operator Hilbert--Schmidt bersifat kompak.
\end{prop}

\begin{proof}[\emph{Petunjuk untuk bukti}] Jika $(e_k)$ suatu basis ortonormal untuk ruang Hilbert tersebut, tinjau
operator $F_n := AP_n$, dengan $P_n$ proyeksi pada rentang linear $\{e_1, \dots, e_n\}$. \ns
\end{proof}

\begin{exam} Operator integral yang didefinisikan dalam contoh~\ref{X_square_int_kernel} merupakan operator Hilbert--Schmidt.
\end{exam}

\begin{exam}
    Operator Volterra yang didefinisikan dalam contoh~\ref{000319} merupakan operator Hilbert--Schmidt.
\end{exam}











\endinput
